\section{Conclusion}
\label{sec:conclusion}

We asked whether a Marchenko--Pastur spectral filter on per-sample
gradients---a mechanism with strong label-noise credentials in
classification---reduces the effect of noise when training on large,
low-SNR financial time-series data. Under a fully pre-registered protocol
with era-purged splits, frozen thresholds, block-bootstrap inference, and
compute-matched tuning, the answer is negative: no evidence of benefit
under the affordable tuning budget, with the filter significantly degrading
out-of-sample performance at the pre-registered operating point on both an
MLP and a GRU, a result that survived independent fresh-seed re-execution
on an independently constructed evaluation split.

The run's equally valuable output is the explanation. For scalar-output
MSE, the filter's engagement eigenspectrum is provably target-independent:
``consensus'' measures feature/Jacobian factor structure, not signal, so
the coherence-amplifier mechanism that worked under label noise has nothing
to select for in this regime. Empirically, MP-eigenselection was
statistically indistinguishable from a norm- and $k$-matched random
subspace projection, and all agreement-style variants hurt similarly---the
harm belongs to the intervention class, not the spectral machinery. Together
these results delimit where coherence-amplifying optimizers can be expected
to apply: the mechanism requires a loss whose per-sample gradient geometry
is target-sensitive, and a regime where coherent target signal is separable
from factor structure---conditions the low-SNR financial regression setting,
by proof and by measurement, does not provide.
